\documentclass[11pt,letterpaper]{article}
% KJO_preamble.tex
% Shared LaTeX preamble for all KJO documentation documents.
% \input this file immediately after \documentclass{}.

% --- Encoding and fonts -------------------------------------------------------
\usepackage[utf8]{inputenc}
\usepackage[T1]{fontenc}
\usepackage{lmodern}
\usepackage{microtype}

% --- Page geometry ------------------------------------------------------------
\usepackage[
  letterpaper,
  top=1in, bottom=1in,
  left=1.25in, right=1in
]{geometry}

% --- Core utilities -----------------------------------------------------------
\usepackage{amsmath}
\usepackage{graphicx}
\usepackage{xcolor}
\usepackage{parskip}        % space between paragraphs; no indent
\usepackage{enumitem}
\usepackage{caption}
\usepackage{subcaption}
\usepackage{float}

% --- Tables -------------------------------------------------------------------
\usepackage{booktabs}
\usepackage{longtable}
\usepackage{array}
\usepackage{multirow}

% --- Code listings ------------------------------------------------------------
\usepackage{listings}
\lstset{
  basicstyle=\ttfamily\small,
  keywordstyle=\bfseries\color{KJOblue},
  commentstyle=\itshape\color{gray},
  stringstyle=\color{KJOgreen!70!black},
  breaklines=true,
  frame=single,
  framesep=4pt,
  xleftmargin=6pt,
  xrightmargin=4pt,
  numbers=left,
  numberstyle=\tiny\color{gray},
  numbersep=8pt,
  language=C++,
}

% --- TikZ / PGF ---------------------------------------------------------------
\usepackage{tikz}
\usetikzlibrary{
  arrows.meta,
  shapes.geometric,
  shapes.misc,
  shapes.symbols,
  positioning,
  fit,
  calc,
  backgrounds,
  decorations.pathreplacing,
  decorations.markings,
  matrix,
}
\usepackage{pgfplots}
\pgfplotsset{compat=1.18}

% --- Headers / footers --------------------------------------------------------
\usepackage{fancyhdr}
\usepackage{lastpage}
\pagestyle{fancy}
\fancyhf{}
\renewcommand{\headrulewidth}{0.4pt}
\renewcommand{\footrulewidth}{0.4pt}
% Per-document: \lhead{}, \rhead{}, \cfoot{} are set in the main .tex file.

% --- Section formatting -------------------------------------------------------
\usepackage{titlesec}
\titleformat{\section}{\large\bfseries\color{KJOblue}}{{\thesection}}{0.8em}{}[\titlerule]
\titleformat{\subsection}{\normalsize\bfseries\color{KJOblue!80!black}}{{\thesubsection}}{0.8em}{}
\titleformat{\subsubsection}{\normalsize\bfseries}{{\thesubsubsection}}{0.8em}{}

% --- Hyperlinks ---------------------------------------------------------------
\usepackage[
  colorlinks=true,
  linkcolor=KJOblue,
  urlcolor=KJOblue!70!black,
  citecolor=KJOblue,
  pdfborder={0 0 0},
]{hyperref}

% --- Color palette -----------------------------------------------------------
\definecolor{KJOblue}   {RGB}{ 30,  90, 160}   % RCU
\definecolor{KJOred}    {RGB}{180,  30,  30}   % EMU
\definecolor{KJOgreen}  {RGB}{ 20, 120,  60}   % GSEMU
\definecolor{KJOorange} {RGB}{200, 110,  20}   % GSE / pressurization
\definecolor{KJOgray}   {RGB}{100, 100, 100}   % Buses / links
\definecolor{KJOpurple} {RGB}{110,  40, 140}   % Shared libraries
\definecolor{KJObg}     {RGB}{248, 248, 252}   % Light background for code

% --- Convenience macros -------------------------------------------------------

% Hardware component name (small caps)
\newcommand{\hw}[1]{\textsc{#1}}

% Software module name (monospace)
\newcommand{\sw}[1]{\texttt{#1}}

% Command name (monospace bold)
\newcommand{\cmd}[1]{\texttt{\textbf{#1}}}

% Pin / signal name
\newcommand{\sig}[1]{\texttt{#1}}

% Note / callout box
\usepackage{tcolorbox}
\tcbuselibrary{skins}
\newtcolorbox{notebox}{
  enhanced,
  colback=KJOblue!5,
  colframe=KJOblue!40,
  fonttitle=\bfseries,
  title=Note,
  left=4pt, right=4pt, top=2pt, bottom=2pt,
}
\newtcolorbox{warningbox}{
  enhanced,
  colback=KJOred!5,
  colframe=KJOred!40,
  fonttitle=\bfseries\color{KJOred},
  title=Warning,
  left=4pt, right=4pt, top=2pt, bottom=2pt,
}

% Title page macro: \KJOtitlepage{title}{subtitle}{version}{date}{description}
\newcommand{\KJOtitlepage}[5]{%
  \begin{titlepage}
    \centering
    \vspace*{2cm}
    {\Huge\bfseries\color{KJOblue} #1 \par}
    \vspace{0.5cm}
    {\Large\color{KJOgray} #2 \par}
    \vspace{1.5cm}
    \rule{0.5\textwidth}{0.4pt}\par
    \vspace{1cm}
    {\large Ken Overton \par}
    \vspace{0.3cm}
    {\normalsize \textit{KJO Liquid Rocket Motor Control System} \par}
    \vspace{1.5cm}
    \begin{tabular}{ll}
      \textbf{Version:} & #3 \\[4pt]
      \textbf{Date:}    & #4 \\
    \end{tabular}
    \vspace{1.5cm}\par
    \begin{minipage}{0.75\textwidth}
      \centering
      \normalsize #5
    \end{minipage}
    \vfill
    {\small\color{KJOgray} \textcopyright\ 2025--2026 Ken Overton.
     All rights reserved. \par}
  \end{titlepage}%
}

\newcommand{\docversion}{0.3}
\newcommand{\docdate}{September 2026}
\lhead{\textsc{KJO Load Cell Management Unit}}
\rhead{\textsc{LCMU Reference} --- v\docversion}
\cfoot{\thepage\ of \pageref{LastPage}}
\hypersetup{pdftitle={KJO LCMU Reference}, pdfauthor={Ken Overton}}

\begin{document}
% KJO_tikz_styles.tex
% Shared TikZ node and edge styles for all KJO block diagrams.
% \input this file inside the document body before any tikzpicture environments.

\tikzset{
  % ── Hardware block (peripheral chip / PCB component) ──────────────────────
  hwblock/.style={
    rectangle, draw=KJOgray, line width=1pt,
    fill=KJOgray!10, rounded corners=4pt,
    minimum width=2.6cm, minimum height=0.8cm,
    font=\small, align=center, inner sep=4pt,
  },
  % Sub-variant: MCU (larger, blue border)
  mcublock/.style={
    hwblock,
    draw=KJOblue, fill=KJOblue!8,
    minimum width=3.0cm, minimum height=1.0cm,
    font=\small\bfseries,
  },
  % Sub-variant: power / battery
  powerblock/.style={
    hwblock,
    draw=KJOorange, fill=KJOorange!8,
  },
  % ── Software module box ────────────────────────────────────────────────────
  swmodule/.style={
    rectangle, draw=KJOblue, line width=0.8pt,
    fill=KJOblue!6, rounded corners=3pt,
    minimum width=2.8cm, minimum height=0.7cm,
    font=\small, align=center, inner sep=3pt,
  },
  % Sub-variant: shared library
  sharedlib/.style={
    swmodule,
    draw=KJOpurple, fill=KJOpurple!6,
    dashed,
  },
  % ── Bus / link annotation ──────────────────────────────────────────────────
  busline/.style={
    draw=KJOgray, line width=1.5pt,
    -{Stealth[length=6pt,width=4pt]},
  },
  bidbusline/.style={
    draw=KJOgray, line width=1.5pt,
    {Stealth[length=6pt,width=4pt]}-{Stealth[length=6pt,width=4pt]},
  },
  radioline/.style={
    draw=KJOblue!60, line width=1.5pt, dashed,
    {Stealth[length=6pt,width=4pt]}-{Stealth[length=6pt,width=4pt]},
  },
  gpioline/.style={
    draw=KJOgreen!60!black, line width=1.2pt,
    -{Stealth[length=5pt,width=4pt]},
  },
  % ── Flowchart shapes ──────────────────────────────────────────────────────
  flowproc/.style={
    rectangle, draw=KJOblue!70, line width=0.8pt,
    fill=KJOblue!8, rounded corners=3pt,
    minimum width=3.2cm, minimum height=0.7cm,
    font=\small, align=center, inner sep=3pt,
  },
  flowdec/.style={
    diamond, draw=KJOorange, line width=0.8pt,
    fill=KJOorange!8,
    minimum width=2.2cm, minimum height=0.9cm,
    font=\small, align=center, inner sep=1pt,
  },
  flowterminal/.style={
    rounded rectangle, draw=KJOblue, line width=1pt,
    fill=KJOblue!12,
    minimum width=2.8cm, minimum height=0.7cm,
    font=\small\bfseries, align=center, inner sep=3pt,
  },
  flowarrow/.style={
    draw=KJOgray, line width=0.9pt,
    -{Stealth[length=5pt,width=4pt]},
  },
  % ── Propellant schematic styles ────────────────────────────────────────────
  % Generic pipe line
  pipe/.style={
    draw=black, line width=1.8pt,
  },
  pipearrow/.style={
    pipe,
    -{Stealth[length=7pt,width=5pt]},
  },
  % Oxidizer (N2O) — blue
  oxpipe/.style={
    draw=KJOblue, line width=1.8pt,
  },
  oxpipearrow/.style={
    oxpipe,
    -{Stealth[length=7pt,width=5pt]},
  },
  % Fuel — red
  fuelpipe/.style={
    draw=KJOred, line width=1.8pt,
  },
  fuelpipearrow/.style={
    fuelpipe,
    -{Stealth[length=7pt,width=5pt]},
  },
  % GSE fill line — green
  gsepipe/.style={
    draw=KJOgreen!70!black, line width=1.8pt,
  },
  gsepipearrow/.style={
    gsepipe,
    -{Stealth[length=7pt,width=5pt]},
  },
  % Pressurization — orange dashed
  presspipe/.style={
    draw=KJOorange, line width=1.4pt, dashed,
  },
  presspipearrow/.style={
    presspipe,
    -{Stealth[length=6pt,width=4pt]},
  },
  % Tank (tall rectangle, rounded corners)
  tank/.style={
    rectangle, draw=black, line width=1.5pt,
    rounded corners=4pt, fill=white,
    minimum width=1.8cm, minimum height=2.4cm,
    font=\small\bfseries, align=center, inner sep=4pt,
  },
  n2otank/.style={
    tank, draw=KJOblue, fill=KJOblue!8,
  },
  fueltank/.style={
    tank, draw=KJOred, fill=KJOred!8,
  },
  srctank/.style={
    tank, draw=KJOgreen!70!black, fill=KJOgreen!6,
  },
  % Valve box (servo-actuated)
  valvebox/.style={
    rectangle, draw=black, line width=1pt,
    fill=white, minimum width=2.2cm, minimum height=0.65cm,
    font=\footnotesize\bfseries, align=center, inner sep=3pt,
  },
  gseval/.style={
    valvebox, draw=KJOgreen!70!black, fill=KJOgreen!8,
  },
  emuval/.style={
    valvebox, draw=KJOred, fill=KJOred!6,
  },
  % Component (generic inline component)
  component/.style={
    rectangle, draw=KJOgray, line width=0.8pt,
    fill=white!96!gray, minimum width=2.2cm, minimum height=0.6cm,
    font=\footnotesize, align=center, inner sep=3pt,
  },
  % Instrument callout circle (PT, TT, TC)
  instrument/.style={
    circle, draw=KJOgray, line width=0.8pt,
    fill=white, minimum size=0.6cm,
    font=\tiny\bfseries, align=center, inner sep=0pt,
  },
  % Sensor lead line
  sensorlead/.style={
    draw=KJOgray, line width=0.7pt,
  },
  % Boundary box (dashed outline for GSE / Rocket regions)
  gseboundary/.style={
    draw=KJOgreen!50, line width=1pt, dashed,
    rounded corners=6pt,
  },
  rocketboundary/.style={
    draw=KJOred!40, line width=1pt,
    rounded corners=6pt,
  },
  % Sequence diagram lifeline
  lifeline/.style={
    draw=KJOgray, line width=0.6pt, dashed,
  },
  % Sequence diagram message arrow
  seqarrow/.style={
    draw=KJOblue, line width=0.9pt,
    -{Stealth[length=5pt,width=4pt]},
  },
  % State node (FSM)
  statenode/.style={
    circle, draw=KJOblue, line width=1pt,
    fill=KJOblue!10,
    minimum size=1.2cm,
    font=\small\bfseries, align=center,
  },
  statetransition/.style={
    draw=KJOgray, line width=0.9pt,
    -{Stealth[length=5pt,width=4pt]},
    font=\footnotesize,
  },
}


\KJOtitlepage%
  {Load Cell Management Unit (LCMU)}%
  {Hardware and Software Reference}%
  {\docversion}{\docdate}%
  {This document describes the hardware, software architecture, CAN bus
   protocol, and load-cell/thrust measurement subsystems of the Load Cell
   Management Unit. The LCMU measures Ox load during oxidizer fill (in
   support of GSEMU's fill-to-target sequence) and measures thrust during
   test-stand burns, logging both to microSD.}

\tableofcontents\newpage

%═══════════════════════════════════════════════════════════════════════════════
\section{Hardware Overview}
%═══════════════════════════════════════════════════════════════════════════════

\subsection{Component Summary}

\begin{center}
\begin{tabularx}{\textwidth}{@{}lp{6.2cm}lY@{}}
\toprule
Component & Part & Interface & Notes \\
\midrule
Microcontroller  & Adafruit Feather RP2040 Adalogger (dual-core Cortex M0+, 133~MHz, 8~MB QSPI flash, 264~KB RAM) & --- & Onboard microSD, STEMMA QT \\
CAN Controller   & Adafruit CAN Bus FeatherWing (MCP2515, TJA1051/3 transceiver) & SPI0 & Link to EMU/GSEMU \\
OLED Display     & Adafruit 128$\times$64 FeatherWing (SH1107), 1.12\textquotedbl, 3 buttons & I\textsuperscript{2}C & Landscape, rotated 180\textdegree \\
Real-Time Clock  & Adafruit DS3231 Precision RTC, STEMMA QT & I\textsuperscript{2}C & Coin-cell (CR1220) backup \\
I\textsuperscript{2}C Multiplexer & Adafruit PCA9548, 8-channel STEMMA QT & I\textsuperscript{2}C & Load-cell ADCs on ports 3, 4, 7 \\
Load Cell ADC (\texttt{$\times$3}) & Adafruit NAU7802, 24-bit, STEMMA QT & I\textsuperscript{2}C (via mux) & One per load cell; fixed address, hence the mux \\
Load Cells (\texttt{$\times$3}) & 0--200~kgf each (0--600~kgf combined), tension/compression, 0.1\% precision & --- & 3~m cable \\
SD Card          & Onboard microSD (dedicated SPI1) & SPI1 & \\
Battery          & 2S LiPo $\to$ external 5~V regulator & ADC (A0) & Resistor-divider sense \\
\bottomrule
\end{tabularx}
\end{center}

\subsection{I\texorpdfstring{\textsuperscript{2}}{2}C Topology and Pin Map}

The DS3231, PCA9548 mux, and all three NAU7802 boards share one STEMMA QT
daisy chain:
\begin{center}
Feather $\longleftrightarrow$ DS3231 (0x68) $\longleftrightarrow$
PCA9548 mux (0x70) $\longleftrightarrow$ ports 3/4/7 $\longleftrightarrow$
3$\times$NAU7802 (0x2A each)
\end{center}
The mux is required because all three NAU7802 boards share the same fixed
I\textsuperscript{2}C address (\texttt{0x2A}) --- only one can be selected
onto the shared bus at a time. Channel select is a direct two-byte
I\textsuperscript{2}C write (\texttt{Wire.beginTransmission(0x70)};
\texttt{Wire.write(1 << channel)}; \texttt{Wire.endTransmission()}), no
dedicated mux library needed.

\begin{center}
\begin{tabularx}{\textwidth}{@{}llY@{}}
\toprule
Pin / Signal & Value & Notes \\
\midrule
\sig{CAN\_CS\_PIN}    & GPIO1 (RX) & \textbf{RX/TX are reversed on this board} versus the Feather M0 convention EMU/GSEMU use (confirmed against Adafruit's pinout page and the earlephilhower core's \texttt{pins\_arduino.h}) \\
\sig{CAN\_INT\_PIN}   & GPIO13 & Not read by firmware (this design polls \sw{parsePacket()} from \sw{loop()}); wired for consistency with EMU/GSEMU and left available for future interrupt-driven use \\
\sig{BUTTON\_A\_PIN}  & GPIO9 & Arms thrust recording (bench testing) \\
\sig{BUTTON\_B\_PIN}  & GPIO6 & Toggle weight recording \\
\sig{BUTTON\_C\_PIN}  & GPIO5 & Tare \\
\sig{BATTERY\_ADC\_PIN} & A0 (GPIO26) & Resistor-divider battery voltage sense, 12-bit ADC \\
\bottomrule
\end{tabularx}
\end{center}

\subsection{Hardware Block Diagram}

\begin{figure}[H]
\centering
\begin{tikzpicture}[node distance=0.9cm and 1.8cm]
  \node[mcublock, minimum width=3.2cm, minimum height=1.1cm] (mcu)
    {Feather RP2040\\Adalogger};
  % I2C chain, left column
  \node[hwblock, left=2.8cm of mcu, yshift= 1.6cm, minimum width=3.0cm] (rtc)  {DS3231 RTC};
  \node[hwblock, below=0.5cm of rtc, minimum width=3.0cm] (mux)  {PCA9548\\I\textsuperscript{2}C Mux};
  \node[hwblock, below=0.4cm of mux, minimum width=2.2cm] (nau1) {NAU7802 (1)};
  \node[hwblock, below=0.35cm of nau1, minimum width=2.2cm] (nau2) {NAU7802 (2)};
  \node[hwblock, below=0.35cm of nau2, minimum width=2.2cm] (nau3) {NAU7802 (3)};
  % SPI, right column
  \node[hwblock, right=2.8cm of mcu, yshift=0.8cm, minimum width=2.8cm] (can) {MCP2515\\CAN FeatherWing};
  \node[hwblock, below=0.6cm of can, minimum width=2.8cm] (sd)  {microSD (onboard)};
  % Below MCU: local UI, then power (offset clear of the OLED node)
  \node[hwblock, below=2.0cm of mcu, minimum width=3.2cm] (oled) {SH1107 OLED\\+ 3 Buttons};
  \node[powerblock, below=1.3cm of oled, xshift=-2.6cm, minimum width=3.0cm] (bat) {2S LiPo + 5V Reg.};
  % I2C bus
  \draw[bidbusline] (mcu.north west) -- ++(-0.9,0) |- (rtc.east) node[near start,above,font=\tiny]{I\textsuperscript{2}C};
  \draw[bidbusline] (mcu.west |- mux) -- (mux.east) node[midway,above,font=\tiny]{I\textsuperscript{2}C};
  \draw[busline] (mux.south) -- (nau1.north);
  \draw[busline] (nau1.south) -- (nau2.north);
  \draw[busline] (nau2.south) -- (nau3.north);
  \draw[bidbusline] (mcu.south) -- (oled.north) node[midway,right,font=\tiny]{I\textsuperscript{2}C};
  % SPI0 (CAN) / SPI1 (SD)
  \draw[bidbusline] (mcu.east |- can) -- (can.west) node[midway,above,font=\tiny]{SPI0};
  \draw[bidbusline] (mcu.east |- sd)  -- (sd.west)  node[midway,above,font=\tiny]{SPI1};
  % Power -- routed around the OLED node, not through it
  \draw[busline,draw=KJOorange] (bat.north) -- ++(0,0.9) -| (mcu.south)
    node[near end,right,font=\tiny]{5V};
\end{tikzpicture}
\caption{LCMU hardware block diagram.}
\end{figure}

%═══════════════════════════════════════════════════════════════════════════════
\section{Software Architecture}
%═══════════════════════════════════════════════════════════════════════════════

\subsection{Module Summary}

\begin{center}
\begin{tabularx}{\textwidth}{@{}llY@{}}
\toprule
Module & Location & Responsibility \\
\midrule
\sw{main.cpp}            & local  & Setup, loop, CAN dispatch, recording modes, buttons, display refresh \\
\sw{KJO\_Load\_Cells}    & local  & PCA9548 mux control, 3$\times$NAU7802 background scan, tare, raw$\to$lbf conversion \\
\sw{KJO\_LCMU\_Display}  & local  & SH1107 OLED idle/run screen \\
\sw{KJO\_LCMU\_Files}    & local  & SD file-sequencing (\sw{Find\_Available\_File()}) \\
\sw{KJO\_Thrust\_Buffer} & local  & Pre/post-trigger RAM ring buffer for high-rate thrust capture \\
\sw{KJO\_CAN\_Command\_Defs} & shared & CAN command codes, packet structs, node IDs \\
\sw{KJO\_Status} / \sw{KJO\_Status\_Display} & shared & Boot-time scrolling status log (\sw{Report\_Status()}) \\
\bottomrule
\end{tabularx}
\end{center}

\subsection{Main Loop Control Flow}

\begin{figure}[H]
\centering
\begin{tikzpicture}[node distance=0.5cm]
  \node[flowterminal] (start)  {loop()};
  \node[flowproc, below=of start]  (scan)  {Cells.scan()};
  \node[flowproc, below=of scan]   (can)   {Check\_CAN()};
  \node[flowproc, below=of can]    (wt)    {Check\_Weight\_Recording()};
  \node[flowproc, below=of wt]     (th)    {Check\_Thrust\_Recording()};
  \node[flowproc, below=of th]     (health) {Check\_Cell\_Health()};
  \node[flowproc, below=of health] (btn)   {Check\_Buttons()};
  \node[flowproc, below=of btn]    (disp)  {Check\_Display()};
  \draw[flowarrow] (start)--(scan);
  \draw[flowarrow] (scan)--(can);
  \draw[flowarrow] (can)--(wt);
  \draw[flowarrow] (wt)--(th);
  \draw[flowarrow] (th)--(health);
  \draw[flowarrow] (health)--(btn);
  \draw[flowarrow] (btn)--(disp);
  \draw[flowarrow] (disp.south) -- ++(0,-0.4) -- ++(4.0,0)
    -- ++(0,6.3) -- (start.east);
\end{tikzpicture}
\caption{LCMU main loop control flow. \sw{Cells.scan()} advances the
background round-robin load-cell read by one step every iteration --
everything else reads back whatever it most recently finished.}
\end{figure}

%═══════════════════════════════════════════════════════════════════════════════
\section{CAN Bus Protocol}
%═══════════════════════════════════════════════════════════════════════════════

Pure addressed command/response: every message is a request or a direct
reply to one, no unsolicited broadcasts. Node IDs: EMU~=~1, GSEMU~=~2,
LCMU~=~3. CAN ID packing (11-bit standard frame):
\texttt{CAN\_ID = (destination << 4) | source}. Baud rate: 1~Mbps.
Frame payload: \texttt{[command(1 byte)][params(up to 7 bytes)]}, defined
in the shared \sw{KJO\_CAN\_Command\_Defs.h} header (separate from the
LoRa-only \sw{KJO\_Command\_Defs.h}, since the two transports serve
different audiences).

\subsection{Command Table}

\begin{center}
\begin{tabularx}{\textwidth}{@{}llp{2.6cm}Y@{}}
\toprule
Command & Sent by & Response & Notes \\
\midrule
\cmd{CAN\_TARE}                   & EMU or GSEMU & ack & Tares all 3 channels (\sw{Cells.tare()}) \\
\cmd{CAN\_GET\_BATTERY\_VOLTAGE}  & EMU          & float volts & \\
\cmd{CAN\_REPORT\_CURRENT\_WEIGHT}& EMU or GSEMU & float lbf & Combined 3-channel total; GSEMU polls this during an active fill \\
\cmd{CAN\_BEGIN\_WEIGHT\_RECORDING} & EMU        & ack/fail & Fails if Thrust mode is active \\
\cmd{CAN\_STOP\_WEIGHT\_RECORDING}  & EMU        & ack      & \\
\cmd{CAN\_BEGIN\_THRUST\_RECORDING} & EMU        & ack/fail & Arms threshold trigger; fails if Weighing mode is active \\
\cmd{CAN\_STOP\_THRUST\_RECORDING}  & EMU        & ack      & Manual override --- forces immediate write-and-close \\
\cmd{CAN\_GET\_HEALTH\_STATUS}      & Any node   & uint16 bitfield & SD/RTC/per-cell fault bits (see below) \\
\cmd{CAN\_PING}                     & Any node   & ack      & Boot-time presence check only \\
\cmd{CAN\_GET\_CELL\_RAW}           & EMU        & float (tare-adjusted raw) & \texttt{param} = channel index (0/1/2); diagnostic readout for RCU\_UNIT\_TEST \\
\bottomrule
\end{tabularx}
\end{center}

\sw{CAN\_GET\_HEALTH\_STATUS}'s response bitfield reuses bits 0--1 for
SD/RTC init failure (shared convention across EMU/GSEMU/LCMU) and bits
8--10 for this board's own 3 load-cell channels
(\texttt{HEALTH\_LCMU\_CELL\_1\_FAIL} through \texttt{\_3\_FAIL}).

\subsection{Fill Sequence (LCMU's Role)}

LCMU does not initiate or track the fill sequence itself --- it only
answers EMU's requests. EMU owns the whole session; GSEMU has no fill
state machine of its own, it only opens/closes the Fill valve servo on
command:
\begin{sloppypar}
\begin{enumerate}[nosep]
  \item Operator commands \cmd{BEGIN\_FILL} (RCU $\to$ EMU).
  \item EMU (fresh start only --- a resumed/paused session skips this,
        since re-taring mid-fill would zero the load cells against mass
        already flowed) sends \cmd{CAN\_TARE} to LCMU.
  \item EMU sends \cmd{CAN\_OPEN\_FILL\_VALVE} to GSEMU.
  \item EMU polls LCMU with \cmd{CAN\_REPORT\_CURRENT\_WEIGHT} every
        200~ms, comparing the tare-relative reading against a per-motor
        target EMU holds locally --- never sent to GSEMU or LCMU.
  \item Once the target is reached, EMU sends \cmd{CAN\_CLOSE\_FILL\_VALVE}
        to GSEMU itself; no round trip through LCMU is needed to
        authorize the close.
\end{enumerate}
\end{sloppypar}

%═══════════════════════════════════════════════════════════════════════════════
\section{Load Cell Subsystem}
%═══════════════════════════════════════════════════════════════════════════════

\subsection{NAU7802 Configuration}

\begin{center}
\begin{tabularx}{\textwidth}{@{}lllY@{}}
\toprule
Constant & Value & Unit & Meaning \\
\midrule
\sw{LOAD\_CELL\_GAIN}            & 128$\times$ & --- & NAU7802 PGA gain \\
\sw{LOAD\_CELL\_EXCITATION\_VOLTS} & 3.0 & V & Bridge excitation, also the ADC reference (ratiometric) \\
\sw{LOAD\_CELL\_SAMPLE\_RATE\_SPS} & 320 & SPS & Fastest available NAU7802 conversion rate \\
\sw{SAMPLE\_AVERAGE\_COUNT}      & 9   & readings & Averaged per reported sample; noise drops as $1/\sqrt{N}$ \\
\sw{TARE\_SAMPLE\_COUNT}         & 10  & readings & Averaged per channel at tare time \\
\sw{LOAD\_CELL\_COUNTS\_PER\_LBF} & 2435.2 & counts/lbf & Estimated (see derivation below); pending known-weight calibration \\
\bottomrule
\end{tabularx}
\end{center}

The excitation voltage is set below the NAU7802's rated 2.4--4.0~V range:
the STEMMA QT bus feeding all three boards runs off the Feather's own
3.3~V rail (not the 5~V the Feather itself is powered from), leaving
little headroom for the chip's internal LDO to regulate cleanly above
3.0~V. Reaching the full rated range for better SNR would require
rewiring the NAU7802 boards' power pin to a 5~V source instead of the
3.3~V STEMMA QT bus --- a hardware change, deferred.

\sw{LOAD\_CELL\_COUNTS\_PER\_LBF} is derived from the load cell's rated
output (1~mV/V at full-scale 200~kgf) and the NAU7802's ratiometric ADC
transfer function:
\[
\text{full-scale ADC code} = \frac{V_{\text{diff}} \times \text{Gain}}{V_{\text{ref}}} \times 2^{23}
= \frac{0.003 \times 128}{3.0} \times 8388608 \approx 1{,}073{,}742 \text{ counts}
\]
\[
\text{full-scale lbf} = 200~\text{kgf} \times 2.20462 \approx 440.92~\text{lbf}
\qquad\Rightarrow\qquad
\text{counts/lbf} = \frac{1{,}073{,}742}{440.92} \approx 2435.2
\]
This is an estimate pending an actual known-weight calibration pass.

\subsection{Background Scan and Tare}

\sw{Load\_Cells::scan()}, called unconditionally every \sw{loop()}
iteration, advances a round-robin background read: one raw attempt per
call for whichever channel is currently up, finalizing that channel's
cached average once \sw{SAMPLE\_AVERAGE\_COUNT} attempts have accumulated
before moving to the next channel. Every other reader
(\sw{readChannels()}, \sw{readWeight()}, \sw{rawCached()}) is instant ---
it reads back the most recently completed batch rather than triggering a
fresh ADC read. A channel reports healthy if at least one reading in its
most recent batch succeeded.

\sw{tareChannel()} averages \sw{TARE\_SAMPLE\_COUNT} raw readings
back-to-back and stores the result as that channel's offset in raw ADC
counts (not lbf) --- this keeps the tare valid independent of
\sw{LOAD\_CELL\_COUNTS\_PER\_LBF}, so a later calibration revision doesn't
invalidate an existing tare. \sw{tare()} tares all 3 channels in turn and
is run automatically as the last step of \sw{setup()}, so the unit starts
up already zeroed, as well as on \cmd{CAN\_TARE} or Button~C.

%═══════════════════════════════════════════════════════════════════════════════
\section{Operating Modes}
%═══════════════════════════════════════════════════════════════════════════════

\sw{current\_mode} (\texttt{MODE\_IDLE} / \texttt{MODE\_WEIGHING} /
\texttt{MODE\_THRUST}) gates both recording behavior and which OLED
buttons are watched. Weighing and Thrust are mutually exclusive --- a
\cmd{CAN\_BEGIN\_THRUST\_RECORDING} while weighing (or vice versa) is
rejected with ack~=~fail.

\subsection{Idle Mode}

Continuously samples all 3 load cells via the background scan, answers
on-demand CAN queries, and watches all three OLED buttons.

\subsection{Weighing Mode}

Opens the next sequenced \texttt{Weight\_N.txt} and appends one
timestamped sample (raw counts and lbf, all 3 channels plus the combined
total) every \sw{WEIGHT\_SAMPLE\_INTERVAL\_MS} (1~s), until
\cmd{CAN\_STOP\_WEIGHT\_RECORDING} or Button~B closes the file. Only
Button~B is watched while weighing; A and C are ignored.

\subsection{Thrust Mode}

\begin{center}
\begin{tabularx}{\textwidth}{@{}lY@{}}
\toprule
Constant & Meaning \\
\midrule
\sw{THRUST\_ARM\_THRESHOLD\_LBF}          & 1.0~lbf --- sustained thrust above this triggers capture \\
\sw{THRUST\_PRE\_TRIGGER\_MS}             & 1000~ms of pre-trigger data buffered continuously while armed \\
\sw{THRUST\_SAMPLE\_INTERVAL\_MS}         & 100~ms (10~Hz) --- same averaged read path as Weighing mode \\
\sw{THRUST\_POST\_TRIGGER\_DURATION\_MS}  & 120{,}000~ms (120~s) --- fixed once triggered, regardless of the thrust profile in between \\
\sw{THRUST\_NEVER\_TRIGGERED\_TIMEOUT\_MS}& 21{,}000~ms (21~s) --- independent safety net so an armed-but-untriggered bench test (no real force applied) still ends on its own \\
\bottomrule
\end{tabularx}
\end{center}

\begin{sloppypar}
\cmd{CAN\_BEGIN\_THRUST\_RECORDING} (or Button~A, bench testing) opens the
next sequenced \texttt{Thrust\_N.txt} and arms \sw{Thrust\_Buffer}, which
continuously buffers the trailing \sw{THRUST\_PRE\_TRIGGER\_MS} of samples
in RAM. Once thrust crosses \sw{THRUST\_ARM\_THRESHOLD\_LBF}, the buffer
freezes and writes that pre-trigger snapshot, then continues sampling into
a second linear (non-overwriting) buffer for the fixed post-trigger
duration. If armed but never triggered, the \emph{never-triggered timeout}
completes the recording anyway --- without it, a bench test with no real
force applied would never end, since the post-trigger completion logic
only starts counting once a trigger has actually occurred. No button is
watched while in Thrust mode; it only ends via one of these two timers or
\cmd{CAN\_STOP\_THRUST\_RECORDING}'s manual override.
\end{sloppypar}

Both timers were revised from the original design during implementation:
real hardware testing found a sustained-low-thrust auto-cutoff (the
original design) risked ending a real burn's recording early on a
transient thrust dip mid-burn, and could never run at all for a bench
test that never crosses the arm threshold in the first place. The
two-independent-timer approach above replaced it.

%═══════════════════════════════════════════════════════════════════════════════
\section{Local UI (OLED + Buttons)}
%═══════════════════════════════════════════════════════════════════════════════

The SH1107 panel is mounted rotated 180\textdegree\ from its natural
orientation. The idle/run screen has 3 fixed lines: the top line shows
the 3 individual load-cell readings in fixed-width fields (left to right
= cells 1, 2, 3); the middle line shows the combined weight/thrust total,
large, bold, and centered in a fixed-width field so the decimal point
never shifts, followed by a small gap and a normal-weight \texttt{lbf}
suffix; the bottom line shows battery voltage right-justified, with the
remaining left-hand space showing a recording/tare status label when
applicable.

\begin{center}
\begin{tabularx}{\textwidth}{@{}lY@{}}
\toprule
Button & Action \\
\midrule
A & \textbf{Idle only.} Arms thrust recording (bench testing) --- one-shot; the recording ends on its own (see Thrust Mode timers above), not via a second press. \\
B & \textbf{Idle:} starts weight recording. \textbf{Weighing:} stops it. Ignored while Thrust recording. \\
C & \textbf{Idle only.} Tare --- see the deliberate on-screen sequence below. \\
\bottomrule
\end{tabularx}
\end{center}

\textbf{Tare sequence} (Button~C, or automatically as the last step of
\sw{setup()}): shows ``Tare'' for 1~s, then for each of the 3 cells in
turn, averages \sw{TARE\_SAMPLE\_COUNT} readings, shows
``\texttt{<N>: <raw average>}'' on the bottom line for 1~s, then moves to
the next cell. Triggered on press (not release) --- the sequence has its
own fixed timing regardless of how long the button is physically held.

\begin{notebox}
An earlier revision let Button~C double as a thrust-stop control, which
made the tare sequence look broken/inconsistent whenever it fired mid-
thrust-test. Button~C is now exclusively the tare button, and no button
is watched at all during Thrust mode.
\end{notebox}

%═══════════════════════════════════════════════════════════════════════════════
\section{SD Logging}
%═══════════════════════════════════════════════════════════════════════════════

Three independent, sequentially-numbered file streams on the SD root
(\sw{Find\_Available\_File()} scans for the lowest unused index):

\begin{center}
\begin{tabularx}{\textwidth}{@{}lY@{}}
\toprule
File & Contents \\
\midrule
\texttt{Log\_N.txt}    & General timestamped event log, same convention as EMU/GSEMU/RCU. One file per power cycle, opened at boot. \\
\texttt{Weight\_N.txt} & Opened on \cmd{CAN\_BEGIN\_WEIGHT\_RECORDING}; one row per 1~Hz sample; closed on stop. \\
\texttt{Thrust\_N.txt} & Opened on \cmd{CAN\_BEGIN\_THRUST\_RECORDING} (arm); written in the two-stage burst described above; closed at auto-cutoff or manual stop. \\
\bottomrule
\end{tabularx}
\end{center}

Each data file (\texttt{Weight\_}/\texttt{Thrust\_}) opens with a header
recording gain, sample rate, excitation voltage, the averaging count, the
current \sw{LOAD\_CELL\_COUNTS\_PER\_LBF} value, and each channel's tare
offset in raw counts --- captured at file-open time rather than referenced
as a constant, since tare offsets change every tare cycle. This lets the
lbf columns be independently reconstructed later even if calibration is
subsequently revised. Both raw counts and converted lbf values are logged
per row for the same reason.

Each mode gets a fresh index the next time it starts, so multiple
weighings/burns in one power cycle don't overwrite each other.

%═══════════════════════════════════════════════════════════════════════════════
\section{Real-Time Clock}
%═══════════════════════════════════════════════════════════════════════════════

DS3231, via \sw{RTClib} (the same library RCU uses). On \sw{begin()}, if
\sw{lostPower()} reports true, the clock is set from compile time
(\texttt{\_\_DATE\_\_}/\texttt{\_\_TIME\_\_}) --- the same auto-heal
pattern EMU's PCF8523 uses.

%═══════════════════════════════════════════════════════════════════════════════
\section{Error Handling}
%═══════════════════════════════════════════════════════════════════════════════

\begin{itemize}[nosep]
  \item \textbf{I\textsuperscript{2}C failures} (DS3231, mux, or any individual NAU7802): logged to \texttt{Log\_N.txt} and shown on the OLED boot status; LCMU continues in a degraded state (computing weight/thrust from whichever channels still respond) rather than halting.
  \item \textbf{CAN failures:} \sw{CAN\_Controller.begin()} failure at boot is logged/shown; LCMU continues local operation (buttons, SD logging) without a working bus link.
  \item \textbf{SD failures:} a failed \sw{SD.begin()} disables file logging but doesn't halt the unit; a failed \sw{SD.open()} for the current mode logs a write-fail message and that recording attempt simply fails (ack~=~fail over CAN) rather than losing buffered data silently.
\end{itemize}

%═══════════════════════════════════════════════════════════════════════════════
\section{Parameters \& Flags}
%═══════════════════════════════════════════════════════════════════════════════

This section collects the tunable operational parameters and feature flags
compiled into this firmware image, together with the source file where each
is defined. It is meant as a single reference point when tracing or changing
a timing, threshold, or calibration value. Pin/channel wiring assignments,
CAN command opcodes, and packet byte layouts are documented separately (the
Hardware Overview pin map and the CAN Bus Protocol sections above) and are
not repeated here. Parameters marked \emph{(shared)} are defined in
\sw{KJO\_\allowbreak Shared\_\allowbreak Libraries} and must agree with the same constant used on
the other board(s) that consume it.

\subsubsection*{Load Cell / NAU7802 Configuration}

\begingroup
\footnotesize
\begin{longtable}{@{}p{3.2cm}p{2.4cm}p{6.2cm}p{2.6cm}@{}}
\toprule
Parameter & Value & Description & Source File \\
\midrule
\texttt{NAU7802\_\allowbreak READ\_\allowbreak TIMEOUT\_\allowbreak MS} & 100~ms & Bounded wait for a channel's ADC conversion to complete before it is marked unhealthy and skipped. & \texttt{KJO\_\allowbreak Load\_\allowbreak Cells.h} \\
\texttt{LOAD\_\allowbreak CELL\_\allowbreak GAIN} & \texttt{NAU7802\_\allowbreak GAIN\_\allowbreak 128} & PGA gain setting applied to each NAU7802 channel at init (128x). & \texttt{KJO\_\allowbreak Load\_\allowbreak Cells.h} \\
\texttt{LOAD\_\allowbreak CELL\_\allowbreak GAIN\_\allowbreak VALUE} & 128 & Human-readable gain value paired with \texttt{LOAD\_\allowbreak CELL\_\allowbreak GAIN}, recorded in data-file headers. & \texttt{KJO\_\allowbreak Load\_\allowbreak Cells.h} \\
\texttt{LOAD\_\allowbreak CELL\_\allowbreak COUNTS\_\allowbreak PER\_\allowbreak LBF} & 2435.2 & Raw-ADC-counts-to-lbf conversion factor; estimated from the load cell's rated output, pending empirical recalibration. & \texttt{KJO\_\allowbreak Load\_\allowbreak Cells.h} \\
\texttt{LOAD\_\allowbreak CELL\_\allowbreak EXCITATION} & \texttt{NAU7802\_\allowbreak 3V0} & NAU7802 LDO excitation/ADC reference voltage setting applied to each channel. & \texttt{KJO\_\allowbreak Load\_\allowbreak Cells.h} \\
\texttt{LOAD\_\allowbreak CELL\_\allowbreak EXCITATION\_\allowbreak VOLTS} & 3.0~V & Human-readable excitation voltage, used in the counts-per-lbf derivation and data-file headers. & \texttt{KJO\_\allowbreak Load\_\allowbreak Cells.h} \\
\texttt{LOAD\_\allowbreak CELL\_\allowbreak SAMPLE\_\allowbreak RATE} & \texttt{NAU7802\_\allowbreak RATE\_\allowbreak 320SPS} & NAU7802 ADC conversion rate, set explicitly since the chip defaults to a much slower 10~SPS. & \texttt{KJO\_\allowbreak Load\_\allowbreak Cells.h} \\
\texttt{LOAD\_\allowbreak CELL\_\allowbreak SAMPLE\_\allowbreak RATE\_\allowbreak SPS} & 320 & Human-readable sample rate paired with \texttt{LOAD\_\allowbreak CELL\_\allowbreak SAMPLE\_\allowbreak RATE}. & \texttt{KJO\_\allowbreak Load\_\allowbreak Cells.h} \\
\end{longtable}
\endgroup

\subsubsection*{Background Scan \& Tare}

\begingroup
\footnotesize
\begin{longtable}{@{}p{3.2cm}p{2.4cm}p{6.2cm}p{2.6cm}@{}}
\toprule
Parameter & Value & Description & Source File \\
\midrule
\texttt{TARE\_\allowbreak SAMPLE\_\allowbreak COUNT} & 10 & Number of raw readings averaged back-to-back to compute each channel's tare offset. & \texttt{KJO\_\allowbreak Load\_\allowbreak Cells.h} \\
\texttt{SAMPLE\_\allowbreak AVERAGE\_\allowbreak COUNT} & 9 & Number of raw readings averaged per reported sample in the background round-robin scan; noise drops as $1/\sqrt{N}$. & \texttt{KJO\_\allowbreak Load\_\allowbreak Cells.h} \\
\texttt{TARE\_\allowbreak DISPLAY\_\allowbreak HOLD\_\allowbreak MS} & 1000~ms & How long each step of the on-screen tare sequence (``Tare,'' then each cell's result) is held on the OLED. & \texttt{main.cpp} \\
\end{longtable}
\endgroup

\subsubsection*{Operating Mode Timing (Weighing \& Thrust)}

\begingroup
\footnotesize
\begin{longtable}{@{}p{3.2cm}p{2.4cm}p{6.2cm}p{2.6cm}@{}}
\toprule
Parameter & Value & Description & Source File \\
\midrule
\texttt{WEIGHT\_\allowbreak SAMPLE\_\allowbreak INTERVAL\_\allowbreak MS} & 1000~ms (1~Hz) & Interval between logged rows while in weight-recording mode. & \texttt{main.cpp} \\
\texttt{THRUST\_\allowbreak ARM\_\allowbreak THRESHOLD\_\allowbreak LBF} & 1.0~lbf & Force level above which a thrust recording is considered triggered. & \texttt{main.cpp} \\
\texttt{THRUST\_\allowbreak ARM\_\allowbreak THRESHOLD\_\allowbreak RAW} & \texttt{THRUST\_\allowbreak ARM\_\allowbreak THRESHOLD\_\allowbreak LBF * LOAD\_\allowbreak CELL\_\allowbreak COUNTS\_\allowbreak PER\_\allowbreak LBF} & Trigger threshold pre-converted to raw ADC counts so the hot sampling loop skips a per-sample lbf division. & \texttt{main.cpp} \\
\texttt{THRUST\_\allowbreak PRE\_\allowbreak TRIGGER\_\allowbreak MS} & 1000~ms (1~s) & Amount of pre-trigger history retained before a thrust trigger. & \texttt{main.cpp} \\
\texttt{THRUST\_\allowbreak POST\_\allowbreak TRIGGER\_\allowbreak DURATION\_\allowbreak MS} & 120000~ms (120~s) & Fixed duration recorded after a thrust trigger, regardless of the thrust profile (bumped from 20~s for slow-starting-motor margin). & \texttt{main.cpp} \\
\texttt{THRUST\_\allowbreak SAMPLE\_\allowbreak INTERVAL\_\allowbreak MS} & 100~ms (10~Hz) & Sampling interval used while in thrust-recording mode. & \texttt{main.cpp} \\
\texttt{THRUST\_\allowbreak NEVER\_\allowbreak TRIGGERED\_\allowbreak TIMEOUT\_\allowbreak MS} & 21000~ms (21~s) & Safety timeout that ends an armed thrust recording if force never crosses the trigger threshold (e.g.\ an unloaded bench test). & \texttt{main.cpp} \\
\texttt{THRUST\_\allowbreak BUFFER\_\allowbreak MAX\_\allowbreak POST\_\allowbreak SAMPLES} & 1210 & Capacity of the post-trigger sample buffer, sized for 120~s at 10~Hz plus a 1~s margin. & \texttt{KJO\_\allowbreak Thrust\_\allowbreak Buffer.h} \\
\texttt{THRUST\_\allowbreak BUFFER\_\allowbreak MAX\_\allowbreak PRE\_\allowbreak SAMPLES} & 10 & Capacity of the pre-trigger ring buffer, sized for $\sim$1~s at 10~Hz. & \texttt{KJO\_\allowbreak Thrust\_\allowbreak Buffer.h} \\
\texttt{THRUST\_\allowbreak BUFFER\_\allowbreak NUM\_\allowbreak CHANNELS} & 3 & Number of per-channel values stored alongside each buffered sample (one per load cell). & \texttt{KJO\_\allowbreak Thrust\_\allowbreak Buffer.h} \\
\end{longtable}
\endgroup

\subsubsection*{OLED/Button UI}

\begingroup
\footnotesize
\begin{longtable}{@{}p{3.2cm}p{2.4cm}p{6.2cm}p{2.6cm}@{}}
\toprule
Parameter & Value & Description & Source File \\
\midrule
\texttt{DISPLAY\_\allowbreak REFRESH\_\allowbreak INTERVAL\_\allowbreak MS} & 500~ms & Minimum interval between OLED idle/run-screen redraws. & \texttt{main.cpp} \\
\texttt{MAX\_\allowbreak LINES} & 8 & Number of most-recent status messages retained and shown in the scrolling boot-status ring buffer. & \texttt{KJO\_\allowbreak Status\_\allowbreak Display.cpp} (shared) \\
\end{longtable}
\endgroup

\subsubsection*{Battery Voltage Monitoring}

\begingroup
\footnotesize
\begin{longtable}{@{}p{3.2cm}p{2.4cm}p{6.2cm}p{2.6cm}@{}}
\toprule
Parameter & Value & Description & Source File \\
\midrule
\texttt{BATTERY\_\allowbreak SCALE\_\allowbreak HIGH} & 7.32 & Numerator of the battery voltage-divider scale factor, from a measured-battery calibration point. & \texttt{main.cpp} \\
\texttt{BATTERY\_\allowbreak SCALE\_\allowbreak LOW} & 2.604 & Denominator of the battery voltage-divider scale factor (ADC pin voltage measured at \texttt{BATTERY\_\allowbreak SCALE\_\allowbreak HIGH}). & \texttt{main.cpp} \\
\texttt{BATTERY\_\allowbreak SCALE} & \texttt{BATTERY\_\allowbreak SCALE\_\allowbreak HIGH / BATTERY\_\allowbreak SCALE\_\allowbreak LOW} ($\sim$2.811) & Combined divider ratio applied to the raw ADC pin voltage to compute actual battery voltage. & \texttt{main.cpp} \\
\texttt{ADC\_\allowbreak REF\_\allowbreak VOLTAGE} & 3.3~V & ADC reference voltage used when converting the raw battery-sense reading to a pin voltage. & \texttt{main.cpp} \\
\texttt{ADC\_\allowbreak MAX\_\allowbreak COUNTS} & 4095.0 & Full-scale ADC count value (12-bit resolution), used to normalize the raw battery-sense reading. & \texttt{main.cpp} \\
\end{longtable}
\endgroup

\subsubsection*{CAN Bus Communication}

\begingroup
\footnotesize
\begin{longtable}{@{}p{3.2cm}p{2.4cm}p{6.2cm}p{2.6cm}@{}}
\toprule
Parameter & Value & Description & Source File \\
\midrule
\texttt{CAN\_\allowbreak BAUDRATE} & 1{,}000{,}000~bps (1~Mbps) & CAN bus bit rate, raised from 250~kbps since the installation's max bus run (16~ft) is well within 1~Mbps margin. & \texttt{KJO\_\allowbreak CAN\_\allowbreak Command\_\allowbreak Defs.h} (shared) \\
\end{longtable}
\endgroup

\subsubsection*{Fill Simulation / Bench-Test Modes (compile-time flags, disabled by default)}

\begingroup
\footnotesize
\begin{longtable}{@{}p{3.2cm}p{2.4cm}p{6.2cm}p{2.6cm}@{}}
\toprule
Parameter & Value & Description & Source File \\
\midrule
\texttt{NO\_\allowbreak LOAD\_\allowbreak CELLS} & undefined (commented out) & Compile-time flag that bypasses real NAU7802 hardware and reports a fixed mid-scale reading (2048 counts) on every channel, for bench-testing without load cells attached. & \texttt{KJO\_\allowbreak Load\_\allowbreak Cells.h} \\
\texttt{LCMU\_\allowbreak FILL\_\allowbreak SIM} & undefined (commented out) & Compile-time flag that makes LCMU fake a ramping fill weight over CAN instead of reading real load cells, for end-to-end EMU/RCU fill-logic testing. & \texttt{main.cpp} \\
\texttt{FILL\_\allowbreak SIM\_\allowbreak RATE\_\allowbreak LBM\_\allowbreak PER\_\allowbreak POLL} & 14.053f / 600.0f ($\sim$0.0234 lbm/poll) & Simulated weight-gain rate per \cmd{CAN\_\allowbreak REPORT\_\allowbreak CURRENT\_\allowbreak WEIGHT} poll, calibrated so the default motor's fill target is reached in $\sim$2 minutes at EMU's poll rate. & \texttt{main.cpp} \\
\texttt{FILL\_\allowbreak SIM\_\allowbreak CEILING\_\allowbreak LBM} & 100.0 & Upper bound on simulated fill weight, set above every real motor's target to prevent runaway growth. & \texttt{main.cpp} \\
\end{longtable}
\endgroup

%═══════════════════════════════════════════════════════════════════════════════
\section{Version History}
%═══════════════════════════════════════════════════════════════════════════════
\begin{center}
\begin{tabularx}{\textwidth}{@{}llY@{}}
\toprule
Version & Date & Notes \\
\midrule
0.1 & September 2026 & Initial release. \\
0.2 & September 2026 & Corrected the Fill Sequence subsection: EMU owns the fill session directly (\sw{CAN\_TARE} and \sw{CAN\_REPORT\_CURRENT\_WEIGHT} sent by EMU itself, target held locally on EMU), not GSEMU as the initial draft stated -- that description matched an earlier, superseded design. \\
0.3 & September 2026 & Add new §10 Parameters \& Flags reference section: tunable timing,
                   threshold, and calibration constants across the load-cell, tare, weighing/thrust,
                   battery, and CAN subsystems, with source-file attribution for each. \\
\bottomrule
\end{tabularx}
\end{center}

\end{document}
